%%============================ Compiler Directives =======================%%
%%                                                                        %%
% !TeX program = lualatex							    	
% !TeX encoding = utf8
% !TeX spellcheck = uk_UA
%%                                                                        %%
%%============================== Клас документа ==========================%%
%%                                                                        %%
\documentclass[14pt]{extarticle}
\IfFileExists{ukrcorr.sty}{\usepackage{ukrcorr}}{}
%%                                                                        %%
%%========================== Мови, шрифти та кодування ===================%%
%%                                                                        %%
\usepackage{fontspec}
\setsansfont{CMU Sans Serif}%{Arial}
\setmainfont{CMU Serif}%{Times New Roman}
\setmonofont{CMU Typewriter Text}%{Consolas}
\defaultfontfeatures{Ligatures={TeX}}
\usepackage[math-style=TeX]{unicode-math}
\usepackage[english, russian, ukrainian]{babel}

%%                                                                        %%
%%============================= Геометрія сторінки =======================%%
%%                                                                        %%
\usepackage[%
	a4paper,%
	footskip=1cm,%
	headsep=0.3cm,% 
	top=2cm, %поле сверху
	bottom=2cm, %поле снизу
	left=2cm, %поле ліворуч
	right=2cm, %поле праворуч
    ]{geometry}                          
%%                                                                        %%        
%%============================== Інтерліньяж  ============================%%
%%                                                                        %%
\renewcommand{\baselinestretch}{1}
%-------------------------  Подавление висячих строк  --------------------%%
\clubpenalty =10000
\widowpenalty=10000
%---------------------------------Інтервали-------------------------------%%
\setlength{\parskip}{0.5ex}%
\setlength{\parindent}{2.5em}%
%%                                                                        %%
%%=========================== Математичні пакети і графіка ===============%%
%%                                                                        %%
\usepackage{amsmath}
\usepackage{graphicx}
\usepackage{floatflt}
%%                                                                        %%
%%========================== Гіперпосилення (href) =======================%%
%%                                                                        %%
\usepackage[%colorlinks=true,
	%urlcolor = blue, %Colour for external hyperlinks
	%linkcolor  = malina, %Colour of internal links
	%citecolor  = green, %Colour of citations
	bookmarks = true,
	bookmarksnumbered=true,
	unicode,
	linktoc = all,
	hypertexnames=false,
	pdftoolbar=false,
	pdfpagelayout=TwoPageRight,
	pdfauthor={Ponomarenko S.M. aka sergiokapone},
	pdfdisplaydoctitle=true,
	pdfencoding=auto
	]%
	{hyperref}
		\makeatletter
	\AtBeginDocument{
	\hypersetup{
		pdfinfo={
		Title={\@title},
		}
	}
	}
	\makeatother 
%%                                                                        %%
%%============================ Заголовок та автори =======================%%
%%                                                                        %%
\title{}
\author{}                                   
%%                                                                        %%
%%========================================================================%%


\begin{document}

Рівняння Ейнштейна
\begin{equation}\label{Ein}
G_{ik} = T_{ik} - \Lambda g_{ik}
\end{equation}

Згортка рівнянь \eqref{Ein}

\begin{equation}\label{EinC}
	 -R = \frac{8\pi G}{c^4} \rho c^2 - 4\Lambda
\end{equation}
\section*{Модель Фрідмана}


Рівняння енергії
\begin{equation}\label{en}
\left({\frac {\dot {a}}{a}}\right)^{2}={\frac {8\pi G\rho }{3}}-\left({\frac {kc^{2}}{a^{2}}}\right)+{\frac {\Lambda c^{2}}{3}}
\end{equation}

Рівняння руху
\begin{equation}\label{mot}
\frac {\ddot {a}}{a}=-{\frac {4\pi G}{3}}\left(\rho +{\frac {3P}{c^{2}}}\right)+{\frac {\Lambda c^{2}}{3}} 
\end{equation}

Рівняння нерозривності
\begin{equation}\label{con}
\frac {d\rho }{dt}=-3\left({\frac {\dot {a}}{a}}\right)\left(\rho +{\frac {P}{c^{2}}}\right).
\end{equation}

\section*{Статичний випадок $\dot {a} = 0$, модель пилевидної матерії $P=0$}

Рівняння енергії
\begin{equation}\label{en1}
0={\frac {8\pi G\rho }{3}}-\left({\frac {kc^{2}}{a^{2}}}\right)+{\frac {\Lambda c^{2}}{3}}
\end{equation}

Рівняння руху
\begin{equation}\label{mot1}
0=-{\frac {4\pi G}{3}}\rho +{\frac {\Lambda c^{2}}{3}} 
\end{equation}

Рівняння нерозривності
\begin{equation}\label{con1}
\frac {d\rho }{dt}=0.
\end{equation}

З рівняння \eqref{en1} 
\begin{equation}\label{rho}
	\rho=\mathrm{const}
\end{equation}

З рівняння \eqref{mot1}
\begin{equation}\label{key}
	\rho = \frac{\Lambda c^2}{4\pi G}
\end{equation}

З цього рівняння густина енергії матерії

\begin{equation}\label{en2}
\rho =\frac{3kc^2}{a^28\pi G} + \frac{\Lambda c^2}{8\pi G}
\end{equation}

Звідки видно, що параметр $k$ ми повинні визначати зі спостережень, іншого виходу нема, рівняння відповіді на це не дають!!!

Якщо наш всесвіт плоский, то $k = 0$, тоді з \eqref{en2}
\begin{equation}\label{en3}
		\rho =\frac{\Lambda c^2}{8\pi G}
\end{equation}

і, логічно, з \eqref{EinC} скалярна кривизна всесвіту $R = 0$.

%Для того, щоб дізнатись $k$ для Всесвіту, треба з якихось дані окремо по $\rho$ і по $\Lambda$.
%
%Знову ж повернемось до статичного Всесвіту. З рівняння \eqref{mot1} та  \eqref{en3} можна визначити $k$, яка в цьому випадку має дорівнювати
%\begin{equation}\label{k}
%	k = \frac23 \Lambda a^2
%\end{equation}

\section*{Як визначити $k$}

Рівняння \eqref{mot1} і \eqref{con1}дають  $\Lambda = \mathrm{const}$. Однак рівняння \eqref{en2} стверджує, що це може бути лише за умови $k = 0$.

\section*{Висновок}

Статичний Всесвіт має бути з плоским!!!

\end{document}


